\ifdefined\gkver\else\def\gkver{student}\fi
\documentclass[\gkver]{gaokaozhenti}
\begin{document}

\chapter{2003年广东广西}

\begin{enumerate}
\item
%% number: 3
%% paperName: 2003年广东广西高考第 3 题 4 分
%% typeId: 1
%% type: 单选题
%% chapter: 气体性质
%% point: 分子力
%% method: 无
%% score: 4
%% degree: 750
%% duplicateId: 0
%% body:
相距很远的两个分子，以一定的初速度相向运动，直到距离最小。在这个过程中，两分子间的分子势能\xzanswer{D}
\fourchoices[answer=D]
{一直增大}
{一直减小}
{先增大，后减小}
{先减小，后增大}
%% answer: D
%% memo:
\memoanswer{
两分子从相距很远到距离最小的过程中，分子间先表现为引力、后表现为斥力：在分子力表现为
引力的阶段，分子力做正功，分子势能减小；当分子间距离减小到平衡位置之后，分子力表现为斥力，
分子力做负功，分子势能增大。
}

\item
%% number: 11
%% paperName: 2003年广东广西高考第 11 题 6 分
%% typeId: 4
%% type: 实验题
%% chapter: 电路
%% point: 多用表的使用
%% method: 无
%% score: 6
%% degree: 850
%% duplicateId: 0
%% body:
下图为一正在测量中的多用电表表盘。
\begin{enumerate}
	\item
	如果是用直流 $10\UV$ 挡测量电压，则读数为\tkanswer[1.6cm]{6.5}$\UV$。
	\item
	如果是用 $\times100$ 挡测量电阻，则读数为\tkanswer[2.2cm]{$8.0\times10^{2}$}$\UO$。
	\item
	如果是用直流 $5\UmA$ 挡测量电流，则读数为\tkanswer[1.6cm]{3.25}$\UmA$。
\end{enumerate}
\onepicture[align=right]{figs/11.png}
%% answer:
\jdanswer{
\begin{enumerate}
	\item
	$6.5\UV$。
	\item
	$8.0\times10^{2}\UO$。
	\item
	$3.25\UmA$。
\end{enumerate}
}
%% memo:
\memoanswer{
（1）直流 $10\UV$ 挡的分度值为 $0.2\UV$，指针位于 6 后第 2.5 格处，读数为 $6.5\UV$。
（2）选用 $\times100$ 挡时，表盘读数为 8.0，故被测电阻为 $8.0\times100\UO=8.0\times10^{2}\UO$。
（3）直流 $5\UmA$ 挡的分度值为 $0.1\UmA$，表盘读数为 $3.25\UmA$。
}

\item
%% number: 18
%% paperName: 2003年广东广西高考第 18 题 13 分
%% typeId: 6
%% type: 计算题
%% chapter: 交流电
%% point: 交流电
%% method: 无
%% score: 13
%% degree: 450
%% duplicateId: 0
%% body:
在图 1 所示区域（图中直角坐标系 $Oxy$ 的 1，3 象限）内有匀强磁场，磁感强度方向垂直于图面向里，大小为 $B$。半径为 $l$，圆心角为 $60^{\circ}$ 的扇形导线框 OPQ 以角速度 $\omega$ 绕 O 点在图面内沿逆时针方向匀速转动，导线框回路电阻为 $R$。
\begin{enumerate}
	\item
	求线框中感应电流的最大值 $I_{0}$ 和交变感应电流的频率 $f$；
	\item
	在图 2 上画出线框转一周的时间内感应电流 $I$ 随时间 $t$ 变化的图象。（规定与图 1 中线框的位置相应的时刻为 $t=0$）
\end{enumerate}
\twopicture[labelA=2003广东广西18a, labelB=2003广东广西18b, align=right]
{figs/18a.png}
{figs/18b.png}
%% answer:
\jdanswer{
\begin{enumerate}
	\item
	$I_{0}=\frac{B\omega l^{2}}{2R}$，$f=\frac{\pi}{\omega}$。
	\item
	如图所示。
	\onepicture[w=6cm]{figs/18_an}
\end{enumerate}
}
%% memo:
\memoanswer{
（1）线框在磁场中转动时，只有一条半径边在磁场边界外切割磁感线，产生的感应电动势为
$E=\frac{1}{2}Bl^{2}\omega$，故感应电流的最大值 $I_{0}=\frac{E}{R}=\frac{B\omega l^{2}}{2R}$。
（2）线框每转过 $\frac{\pi}{3}$ 的圆心角，切割磁感线的半径边更换一次，感应电流方向改变一次，
故其随时间变化的图象如答案图所示。
}

\end{enumerate}

\end{document}
